% 对应题号 M1-58，复用已核对的本题几何图。
% 修订：启用 xeCJK 自动换行；保留原题参数与物理关系。
\documentclass[tikz,border=14pt]{standalone}
\usepackage{fontspec}
\usepackage{xeCJK}
\usepackage{amsmath,amssymb}
\usepackage{xcolor}
\usepackage{tikz}
\usetikzlibrary{arrows.meta,calc,patterns,positioning,decorations.pathmorphing}
\IfFontExistsTF{Heiti SC}{\IfFontExistsTF{Heiti SC}{\setmainfont{Heiti SC}}{\setmainfont{FandolHei-Regular}}\IfFontExistsTF{Heiti SC}{\setsansfont{Heiti SC}}{\setsansfont{TeX Gyre Heros}}}{\IfFontExistsTF{Songti SC}{\setmainfont{Songti SC}}{\setmainfont{FandolSong-Regular}}\IfFontExistsTF{Songti SC}{\setsansfont{Songti SC}}{\setsansfont{FandolSong-Regular}}}
\definecolor{bluefield}{RGB}{24,76,255}
\definecolor{redfield}{RGB}{232,49,35}
\definecolor{orangefield}{RGB}{224,116,10}
\definecolor{greenfield}{RGB}{18,132,86}
\definecolor{purplefield}{RGB}{112,72,190}
\definecolor{linegray}{RGB}{60,68,78}
\definecolor{softblue}{RGB}{238,243,255}
\definecolor{softgreen}{RGB}{238,249,244}
\definecolor{softorange}{RGB}{255,247,235}
\definecolor{softpurple}{RGB}{247,242,255}
\tikzset{
  arrow/.style={-{Latex[length=3.1mm,width=2.1mm]},line width=1.1pt},
  thinarr/.style={-{Latex[length=2.5mm,width=1.7mm]},line width=.85pt},
  axis/.style={-{Latex[length=2.5mm,width=1.7mm]},line width=.8pt,draw=linegray},
  guide/.style={draw=linegray!55,dashed,line width=.65pt},
  fieldline/.style={-{Latex[length=2.5mm,width=1.7mm]},line width=1pt,draw=bluefield},
  box/.style={draw=linegray,rounded corners=3pt,fill=white,align=center,inner sep=6pt,line width=.75pt},
  bluebox/.style={box,draw=bluefield,fill=softblue},
  greenbox/.style={box,draw=greenfield,fill=softgreen},
  orangebox/.style={box,draw=orangefield,fill=softorange},
  purplebox/.style={box,draw=purplefield,fill=softpurple}
}
\IfFontExistsTF{Heiti SC}{\setCJKmainfont[AutoFakeBold=1.5]{Heiti SC}}{\setCJKmainfont[AutoFakeBold=1.5]{FandolHei-Regular}}
\IfFontExistsTF{Heiti SC}{\setCJKsansfont{Heiti SC}}{\setCJKsansfont{FandolHei-Regular}}
\begin{document}

\begin{tikzpicture}[x=1cm,y=1cm]
\path[use as bounding box] (-10,-4.6) rectangle (10,4.6);
\fill[white] (-10,-4.6) rectangle (10,4.6);
\node[font=\Large\bfseries,text width=18.4cm,align=center] at (0,4.05) {例：挖去奇点后，外边界通量由内边界补偿};
\begin{scope}[shift={(-4.8,-.25)}]
  \fill[softblue] (0,0) circle(2.8); \fill[white] (0,0) circle(.65);
  \draw[linegray] (0,0) circle(2.8); \draw[linegray,dashed] (0,0) circle(.65);
  \fill[redfield] (0,0) circle(.09) node[below right] {奇点};
  \foreach \a in {0,45,...,315}{\draw[arrow,bluefield] ({1.15*cos(\a)},{1.15*sin(\a)})--({2.25*cos(\a)},{2.25*sin(\a)});}
  \node at (0,-3.35) {$V_\varepsilon=V\setminus B_\varepsilon$};
\end{scope}
\node[greenbox,text width=7.4cm] at (4.8,1.05) {$\nabla\cdot\mathbf F=0\quad(r>0)$\\$\displaystyle \iint_{S_R}\mathbf F\cdot d\mathbf S+\iint_{S_\varepsilon}^{\text{区域外法向}}\mathbf F\cdot d\mathbf S=0$};
\node[orangebox,text width=7.4cm] at (4.8,-1.55) {令 $\varepsilon\to0$ 后，内边界通量不会消失，而是成为点源的 $\delta$ 贡献。};
\end{tikzpicture}
\end{document}
